\section{Finite dimensionality of the original automorphic row spaces}
\label{sec:row-finiteness-v47}

This chapter proves the full original \texttt{RowsFinDim} assertion.
The test function need not be an idempotent, and its archimedean unit
vector need not be smooth or compact-finite.  There is no gap, rank,
weight, cuspidality or ordinary hypothesis beyond the literal statement.
We use classical Harish--Chandra finiteness in its precise published
scope; the transfer to the original quotient, level and operators is
proved here.  Source commit:
\texttt{7cbf7a0c32fca1badf05931879a37c3f8a2b03b4}.

Fix an original \(I:\texttt{AutInst}\).  Its group is the unitary
group of the supplied nonsingular Hermitian matrix over its CM
extension \(E/F\).  Write
\[
 G_A=G_\infty\times G_f,
 \quad \Gamma=G(F),
 \quad \mathcal H=L^2(\Gamma\backslash G_A,\nu).
 \tag{RFD.1}
\]
Here \(\nu\) is the quotient measure from the supplied Haar measure
\(\mu_f\), the original infinity Haar measure, and the supplied
fundamental domain.  No probability normalization of \(\nu\) is made.
The operator \(R(g)\) is the original right regular representation.

% BEGIN RFD1
\paragraph{RFD1. The exact classical finite-multiplicity input.}
\label{claim:RFD1}
Let \(G_\infty\) be a product of real unitary groups and let
\(\Delta\) be an arithmetic subgroup of a rational reductive group
whose real group is \(G_\infty\).  For every irreducible continuous
unitary representation \(\pi\), the vector space of bounded
intertwiners
\[
 \operatorname{Hom}^{\mathrm{bdd}}_{G_\infty}
       (\pi,L^2(\Delta\backslash G_\infty))
 \tag{RFD.2}
\]
is finite dimensional, with any positive Haar normalization.

\emph{Published input and hypothesis match.}
The inputs are Harish--Chandra admissibility of irreducible unitary
representations and finiteness of automorphic forms of a fixed finite
compact type and fixed finite-codimensional central ideal.  A primary
exposition is A.~Borel, \emph{Automorphic forms on reductive groups},
\S6.5, \S8.4 and \S10, printed pages 17, 23 and 27--28; in the latter section
the subparagraphs are printed as 8.5--8.6.  The same text explicitly
identifies compact-finite vectors in an irreducible \(L^2\) submodule
as automorphic forms, and derives finite multiplicity.  It is available at
\url{https://homepages.math.uic.edu/~rtakloo/papers/borel/borel3.pdf}.
These are named classical foundations, not new proofs of the
Harish--Chandra theorems.

For clarity, the deduction that will be used is as follows.  Fix a
nonzero smooth compact-finite vector \(v\) in \(\pi\); its existence
and its infinitesimal character follow from admissibility.  Let
\(\xi\) be a finite sum of compact-type idempotents fixing \(v\),
and let \(J\) be the kernel of that infinitesimal character.
For every bounded intertwiner \(A\), the vector \(A(v)\) has the
same compact-type and central-annihilator restrictions.  Its
automorphic representative is therefore in the one finite-dimensional
space \(\mathcal A(\Delta,J,\xi)\).  The map \(A\mapsto A(v)\)
is injective: the closed span of the translates of \(v\) is all of
\(\pi\), and a bounded map vanishing there is zero.  This proves
(RFD.2).  If \(A\ne0\), Schur's lemma gives \(A^*A=c\,1\),
\(c>0\), so its image is an actual closed irreducible submodule of
\(L^2\); the cited assertion about \(L^2\) vectors applies.
This deduction does not assume that the full regular representation
is discrete or that its residual spectrum vanishes.

The groups in question are connected reductive real matrix groups
with compact center, with arithmetic \(\Delta\).  These are precisely
the classical hypotheses.  Equivalently one can restrict scalars to
\(\mathbf Q\); there is no split central factor needing a separate
central-character quotient.  Multiplying Haar measure by a positive
constant gives an isometric rescaling of the Hilbert space and leaves
the dimension in (RFD.2) unchanged.  The trivial rank-zero group is
handled directly and has a one-dimensional regular space.
% END RFD1

% BEGIN RFD2
\paragraph{RFD2. One compact-open level for the complete original test function.}
\label{claim:RFD2}
Suppose \(e:G_f\to\mathbf C\) is locally constant and
\(e(x)\ne0\Rightarrow x\in I.Kset\).  There is a compact open
subgroup \(K_0\subset I.Kset\) such that
\[
             e(kx)=e(x)=e(xk)\qquad(k\in K_0, x\in G_f).
 \tag{RFD.3}
\]
In particular \(e\) is bounded and integrable.

\emph{Proof.}
The stabilizer \(I.Kset\) of the original full integral lattice is
compact open.  At each finite place choose an integral lattice basis;
its stabilizer is a closed subgroup of the compact matrix group,
and a sufficiently small principal congruence subgroup lies in it.
The restricted product has the same property.  This is the compact
open subgroup used in the original adelic topology, as in FV4.

Each point of \(I.Kset\) has an open right coset on which \(e\)
is constant.  Choose a finite such cover by compactness.  Intersect
the finitely many open subgroups, and then take its core in
\(I.Kset\).  The core is still open: an open subgroup of a compact
group has finite index, so there are only finitely many conjugates.
Call the resulting open normal subgroup \(K_0\).  Every chosen
coset is a union of \(K_0\)-cosets; hence \(e\) is constant on
right \(K_0\)-cosets in \(I.Kset\).  Normality gives left
invariance there.  Outside \(I.Kset\), multiplication by an element
of \(K_0\) stays outside, and both sides are zero.  The finite
coset cover bounds the values of \(e\), and its support lies in
the compact set \(I.Kset\), proving integrability.  The level is
chosen from \(e\) alone, before choosing any archimedean vector.
% END RFD2

% BEGIN RFD3
\paragraph{RFD3. The actual fixed-level quotient has finite real multiplicities.}
\label{claim:RFD3}
For a compact open \(K_0\subset G_f\), put
\(\mathcal H^{K_0}=\{F:R(1,k)F=F\ (k\in K_0)\}\).
Every irreducible unitary representation of \(G_\infty\) has
finite bounded-intertwiner multiplicity in \(\mathcal H^{K_0}\).

\emph{Proof.}
The coordinate argument in FV1--FV4 applies to the literal
\(I.E/I.F,I.J,I.L\): the unitary equation defines a connected
reductive group, its geometric form is \(\mathrm{GL}_n\),
and conjugation negates determinant characters.  Thus it has no
nontrivial rational character, and its real center is compact.
FV5--FV10 consequently apply to these same coordinates and give
finitely many representatives \(s\in G_f\) for
\(G(F)\backslash G_f/K_0\).  For each put
\[
 \Delta_s=\{\gamma_\infty:\gamma\in G(F),\
                         \gamma_f\in sK_0s^{-1}\}.
 \tag{RFD.4}
\]
The arithmeticity proof in FV10 uses the inverse image of
\(sK_0s^{-1}\), its finite-index intersection with an integral
stabilizer, and restriction of scalars.  It therefore proves
arithmeticity for these very subgroups; no strong approximation
or single-class assertion is being used.

A right \(K_0\)-invariant function is determined by the finitely
many functions \(F_s(g)=F(g,s)\).  They are left
\(\Delta_s\)-invariant.  Conversely these functions reconstruct
the original function using \(g_f=\gamma_fs k\); independence
of the representation follows exactly from (RFD.4).
Unfolding the original Haar measure gives the direct sum of the
\(L^2(\Delta_s\backslash G_\infty)\) norms with positive
constant factors.  These are the fiber volumes of \(K_0\) in
the original product Haar convention.  In particular the map is
a Hilbert equivalence after multiplying each component by the
square root of its factor.  Null-set identifications and the
choice of the original fundamental domain are legitimate by
RQC2--RQC3.  The equivalence commutes with the original right
\(G_\infty\)-action.

Apply RFD1 to each of the finitely many arithmetic components.
The space of intertwiners into their finite Hilbert direct sum
is the direct sum of the corresponding finite-dimensional spaces.
Transporting back proves the assertion with the original measure
and the original regular action.  No change of those data is
made in the conclusion.
% END RFD3

% BEGIN RFD4
\paragraph{RFD4. The original infinity tensor representation and coefficient.}
\label{claim:RFD4}
Let \((H_v,\sigma_v,d_v,w_v)\) be the complete original infinity
data in \texttt{RowsFinDim}.  On their Hilbert tensor product set
\[
 H_\sigma=\widehat\bigotimes_{v\mid\infty}H_v,
 \quad \sigma=\widehat\bigotimes_v\sigma_v,
 \quad w=\bigotimes_vw_v,
 \quad d=\prod_vd_v.
 \tag{RFD.5}
\]
This is an irreducible strongly continuous unitary representation,
\(\|w\|=1\), and, with inner products linear in the second variable,
\[
 m(g)=\langle w,\sigma(g)w\rangle,
 \qquad \phi(g)=d\,\overline{m(g)}
       =\texttt{phInf}(I,\sigma_v,w_v,d_v)(g).
 \tag{RFD.6}
\]
The function \(\phi\) is integrable for the original infinity measure.

\emph{Proof.}
There are finitely many real places.  Start with the algebraic
tensor product, give elementary tensors the product inner product,
and complete.  The tensor operators are isometries and extend to
that completion.  Continuity follows on elementary tensors from
the finite product, and then on all vectors by isometric
approximation.  The elementary tensors are dense.  Each factor
is irreducible by the original hypothesis; its bounded commutant
is scalar.  The bicommutant theorem therefore makes the strong
closure of its operator algebra all bounded operators.  Acting
one factor at a time shows that the commutant of the product
representation is scalar, so the product representation is
irreducible.  The empty tensor convention is \(\mathbf C\).

The inner-product identity on tensors gives (RFD.6), the precise
product used by the source.  Tonelli applied to absolute values
gives
\(\int|\phi|=\prod_v\int|\phi_v|<\infty\).
The original four-vector formal-degree identities tensor by
Fubini.  No compact-finiteness of the chosen \(w_v\) is used.
% END RFD4

% BEGIN RFD5
\paragraph{RFD5. Schur compression on every unitary Hilbert representation.}
\label{claim:RFD5}
Let \(U\) be any strongly continuous unitary representation of
\(G_\infty\), and let \(T=U(\phi)\), with \(\phi\) from RFD4.
Then
\[
 T^*=T,\qquad T^2=T,\qquad T U(g)T=m(g)T.
 \tag{RFD.7}
\]
These conclusions also hold for an operator given by the original
weak \texttt{IsOpInt} identity.

\emph{Proof.}
The original four-vector Schur identity, applied one place at a
time and then multiplied, gives the coefficient convolution laws
\[
 \phi(g^{-1})=\overline{\phi(g)},\qquad
 \int\phi(h)\phi(x^{-1}h^{-1}y)\,dh=m(x)\phi(y).
 \tag{RFD.8}
\]
To check the second equation, substitute
\(\phi(h)=d\overline{\langle w,\sigma(h)w\rangle}\), move
\(\sigma(x)\) and \(\sigma(y)\) across the inner product by
unitarity, and use the four-vector identity; its factor \(1/d\)
cancels one of the two factors \(d\).  The remaining coefficient
is \(m(x)\), with the displayed order of \(x^{-1}h^{-1}y\).
Putting \(x=1\) gives \(\phi*\phi=\phi\).

Since \(\phi\in L^1\), its integrated operator exists and has
norm at most \(\|\phi\|_1\).  The bounds for products of two
integrands are multiples of \(|\phi(h)\phi(k)|\), so Fubini
justifies operator multiplication through weak matrix coefficients.
Inversion and (RFD.8) prove all of (RFD.7).  A weakly specified
operator agrees with this integrated operator because its inner
product with every vector agrees.  This is also the direct
Schur-compression computation underlying ARO16; here its stated
formal-degree hypotheses, rather than \texttt{GapHyp}, are used.
% END RFD5

% BEGIN RFD6
\paragraph{RFD6. Every projector vector gives a bounded intertwiner.}
\label{claim:RFD6}
For every \(v\in\operatorname{ran}T\), there is a unique bounded
intertwiner \(A_v:H_\sigma\to H_U\) satisfying
\[
 A_v(\sigma(g)w)=U(g)v,
 \qquad A_v(w)=v,
 \qquad \|A_v\|=\|v\|.
 \tag{RFD.9}
\]
The map \(v\mapsto A_v\) is linear and injective.

\emph{Proof.}
For \(v,v'\in\operatorname{ran}T\), self-adjointness and
compression give
\[
 \langle U(g)v,U(h)v'\rangle
   =m(g^{-1}h)\langle v,v'\rangle.
 \tag{RFD.10}
\]
The same coefficient \(m(g^{-1}h)\) is
\(\langle\sigma(g)w,\sigma(h)w\rangle\).
For a finite linear combination, (RFD.10) therefore shows
\[
 \left\|\sum_j a_jU(g_j)v\right\|^2
 =\|v\|^2\left\|\sum_j a_j\sigma(g_j)w\right\|^2.
 \tag{RFD.11}
\]
Thus the prescribed map is well-defined on the algebraic orbit
span and is bounded.  That span is dense by irreducibility and
\(w\ne0\); completeness extends it uniquely.  The formula on
the orbit span proves intertwining, linearity in \(v\), and
uniqueness.  Evaluation at \(w\) proves injectivity and the
operator-norm equality.  This gives a direct argument even when
the ambient regular representation has continuous spectrum.
% END RFD6

% BEGIN RFD7
\paragraph{RFD7. The fixed-level archimedean projector has finite rank.}
\label{claim:RFD7}
For every compact open \(K_0\subset G_f\), the restriction of the
original \(T=R(\phi)\) to \(\mathcal H^{K_0}\) has finite-dimensional
range.

\emph{Proof.}
Finite and infinite right translations commute, so their weak
integrals commute, and \(T\) preserves \(\mathcal H^{K_0}\).
Apply RFD6 to the infinity representation on this closed Hilbert
subspace.  It injects \(\operatorname{ran}(T|\mathcal H^{K_0})\)
linearly into
\[
 \operatorname{Hom}^{\mathrm{bdd}}_{G_\infty}
       (H_\sigma,\mathcal H^{K_0}),
\]
which is finite dimensional by RFD3--RFD4.  This proves the
assertion, with the actual original \(w\); the auxiliary
compact-finite vector in RFD1 was used only to bound the space
of intertwiners and does not replace \(w\).
% END RFD7

% BEGIN RFD8
\paragraph{RFD8. Full original RowsFinDim, including arbitrary test functions.}
\label{claim:RFD8}
For every datum in the full original \texttt{RowsFinDim} statement,
\[
 \dim_{\mathbf C}\operatorname{ran}(R_eT)<\infty.
 \tag{RFD.12}
\]
Consequently the literal assertion
\(\texttt{Formalism.L2.SG.RowsDef.RowsFinDim}\) holds at the
classical foundations recorded in RFD1 and FV5--FV10.

\emph{Proof.}
Choose \(K_0\) by RFD2 and let
\(P_0=\mu_f(K_0)^{-1}\int_{K_0}R(1,k)\,dk\).
It is the orthogonal projection onto \(\mathcal H^{K_0}\).
Both sides of (RFD.3) and Haar change of variables give
\[
                 P_0R_e=R_e=R_eP_0.
 \tag{RFD.13}
\]
The finite and infinity actions commute.  Fubini is legitimate
because \(e,\phi\in L^1\); hence \(R_eT=TR_e\).
For any \(F\in\mathcal H\), the vector \(R_eF\) is in
\(\mathcal H^{K_0}\), so
\[
 R_eTF=T(R_eF)\in\operatorname{ran}(T|\mathcal H^{K_0}).
\]
RFD7 proves this containing space finite dimensional.  The
original \texttt{IsRConvF} and \texttt{IsOpInt} identities
uniquely identify \(R_e,T\), so this is precisely the original
\texttt{HH} range on the supplied \(L^2\) space.  Neither
\(e*e=e\) nor \(e^*=e\) has been used; neither is a hypothesis
of \texttt{RowsFinDim}.  All coefficient-field, signature, Haar
and fundamental-domain quantifiers remain those of the source.
% END RFD8

% BEGIN RFD9
\paragraph{RFD9. The finite-dimensionality input is removed from the row constructions.}
\label{claim:RFD9}
On every supplied original adelic \(\texttt{SG.HidaSetting}\),
the finite-multiplicity hypothesis retained in ARO18--ARO30 is
now supplied by RFD8.  Their actual row construction, including
\texttt{RowsAdd}, \texttt{RowsGood}, \texttt{RowsLab} and one
common \texttt{AllRows}, consequently requires only their other
stated inputs.

\emph{Proof.}
The specialization in ARO18 supplies the original infinity
formal-degree and integrability hypotheses from the recorded
\texttt{GapHyp}, and its event function is locally constant
with support in \(I.Kset\).  RFD8 therefore proves the exact
finite-dimensionality assertion for that event.  The original
evaluation bijection in ARO17 transfers it to the bounded
intertwiner space.  ARO19--ARO24 then construct the rows with
the same measures, projectors, labels and ties.  ARO25--ARO30
use these same rows and do not require a further finiteness
input.  Their original \texttt{AutTie}, \texttt{RowDictTie},
\texttt{LabTie}, arithmetic counting and common-family premises
are unchanged.  Thus this removes \texttt{RowsFinDim} as a
producer task without closing any of those distinct premises.
% END RFD9
